\documentclass[10pt,a4paper]{article}
\usepackage[margin=2.6cm]{geometry}
\usepackage[T1]{fontenc}
\usepackage{amsmath,amssymb,amsthm}
\IfFileExists{newtxtext.sty}{\usepackage{newtxtext,newtxmath}}{\usepackage{txfonts}}
\usepackage[scaled=.95]{helvet}
\usepackage{mathpartir,microtype,booktabs,tabularx,listings,xcolor,titlesec,enumitem}
\usepackage[most]{tcolorbox}
% Palette shared with Tolerant frames.
\definecolor{accent}{HTML}{1F5F8B}
\definecolor{chgc}{HTML}{0B7A68}
\definecolor{rulec}{HTML}{8A1C4A}
\definecolor{codec}{HTML}{24292F}
\definecolor{ckc}{HTML}{0550AE}
\definecolor{csc}{HTML}{116329}
\definecolor{cnc}{HTML}{953800}
\definecolor{ccmc}{HTML}{6E7781}
\definecolor{leanc}{HTML}{2F6F5E}
\definecolor{defbg}{HTML}{EEF6F1}
\definecolor{exbg}{HTML}{F6F4EE}
\usepackage[colorlinks,linkcolor=accent,citecolor=accent,urlcolor=accent]{hyperref}
\titleformat*{\section}{\Large\bfseries\sffamily\color{accent}}
\titleformat*{\subsection}{\large\bfseries\sffamily\color{accent}}
\titleformat*{\paragraph}{\bfseries\sffamily\color{accent}}
\newcommand{\ruleName}[1]{\textcolor{rulec}{\textsc{#1}}}
\lstset{basicstyle=\small\ttfamily,columns=fullflexible,keepspaces=true,
  breaklines=true,showstringspaces=false,upquote=true,xleftmargin=1em,
  frame=l,rulecolor=\color{accent},framesep=6pt,backgroundcolor=\color{exbg},
  keywordstyle=\color{ckc}\bfseries,
  commentstyle=\color{ccmc},stringstyle=\color{csc}}
\lstdefinelanguage{TypeScript}{
  morekeywords={type,export,function,const,let,while,for,if,else,throw,new,return,
    readonly,number,string,void,boolean,extends,true,false,null,undefined,
    Record,SafeInt,Dense,Fresh,Footprint,Modifies,AtMost,Effects,Join},
  sensitive=true,morecomment=[l]{//},morecomment=[s]{/*}{*/},
  morestring=[b]",morestring=[b]',
}
\theoremstyle{definition}
\newtheorem{definition}{Definition}
\newtheorem{example}[definition]{Example}
\newtheorem{theorem}{Theorem}
\newtheorem{proposition}[theorem]{Proposition}
\newtheorem{conjecture}{Conjecture}
\newcommand{\paperbox}[2]{\tcolorboxenvironment{#1}{enhanced,breakable,
  colback=#2,colframe=#2,boxrule=0pt,frame hidden,sharp corners,
  borderline west={2.5pt}{0pt}{accent},left=9pt,right=6pt,top=4pt,bottom=4pt,
  before skip=10pt,after skip=10pt}}
\paperbox{definition}{defbg}
\paperbox{theorem}{defbg}
\paperbox{proposition}{defbg}
\paperbox{conjecture}{exbg}
\paperbox{example}{exbg}

\newcommand{\change}[1]{\textcolor{chgc}{\mathsf{#1}}}
\newcommand{\keep}{\change{keep}}
\newcommand{\init}{\change{init}}
\newcommand{\scalar}{\change{scalar}}
\newcommand{\extend}{\change{extend}}
\newcommand{\grow}{\change{grow}}
\newcommand{\any}{\change{any}}
\newcommand{\kw}[1]{\mathsf{#1}}
\newcommand{\val}{\kw{values}}
\newcommand{\st}{\kw{spines}}
\newcommand{\vw}{\kw{views}}
\newcommand{\LF}{\mathit{Loc}_{\kw{field}}}
\newcommand{\code}[1]{\textcolor{codec}{\texttt{#1}}}
\ExplSyntaxOn
\NewDocumentCommand{\lean}{m}{
  \textcolor{leanc}{\texttt{\small
    \tl_set:Nn \l_tmpa_tl { #1 }
    \tl_replace_all:Nnn \l_tmpa_tl { . } { .\allowbreak }
    \tl_replace_all:Nnn \l_tmpa_tl { \_ } { \_\allowbreak }
    \tl_use:N \l_tmpa_tl
  }}
}
\ExplSyntaxOff
\newcommand{\status}[1]{\textcolor{leanc}{\textsf{\small[#1]}}}
\newcommand{\sem}[1]{[\![#1]\!]}
\newcommand{\tr}[1]{\xrightarrow{#1}}
\newcommand{\sep}{\mathrel{\#}}
\newcommand{\len}[1]{|#1|}
\newcommand{\Lang}{\mathcal{L}}
\newcommand{\Reg}{\mathit{Reg}}
\newcommand{\Eff}{\mathit{Eff}}
\newcommand{\SemEff}{\mathit{SemEff}}
\newcommand{\EExp}{\mathit{EExp}}
\newcommand{\Var}{\mathit{Var}}
\newcommand{\EVar}{\mathit{EVar}}
\newcommand{\Loc}{\mathit{Loc}}
\newcommand{\Val}{\mathit{Val}}
\newcommand{\roots}{\mathrm{roots}}
\newcommand{\reach}[2]{#1.{\ast\ast}_{#2}}
\newcommand{\modstate}[1]{#1.{@\kw{state}}}
\newcommand{\hoare}[3]{\{#1\}\;#2\;\{#3\}}
\newcommand{\Same}{\mathrm{Same}}
\newcommand{\Anch}{\mathrm{Anchored}}

\title{\sffamily\bfseries\color{accent}An Effect Language for TScert\\[2pt]
  \large\normalfont\color{black}Relational Effects, Bounded Effect Polymorphism, and Tolerant Framing}
\author{TScert Working Paper}
\date{27 November 2026}

\emergencystretch=3em
\begin{document}
\maketitle

\begin{abstract}\noindent
An effect bounds how a computation step may modify the heap. This paper
presents a unified effect language for TScert. Each effect denotes a
binary relation between pre- and post-step heaps: its entries assign
change bounds from a lattice $\Delta$ to heap regions, and an explicit
$\kw{else}$ entry bounds every uncovered location. Under this
representation, both the write effect of a command ($\kw{else}:\keep$)
and the preservation support of a fact ($\kw{else}:\any$) are effects, so
framing, subeffecting, and callback bound checking all reduce to one
inclusion relation. To verify higher-order functions such as collection
and text walkers, the language adds capture-rooted regions, type-driven
declared reach ($\reach{\pi}{T}$), module-private confined state
($\modstate{M}$), and bounded effect polymorphism via nine rules
(\ruleName{Bound}, \ruleName{Anch}, \ruleName{Meet}, \ruleName{Inst},
\ruleName{Confined}, \ruleName{Callback-Inv}, \ruleName{Fields-Disj},
\ruleName{Plain-Cells-Disj}, \ruleName{Elem-Len-Extend}). TScert's
verified footprint fragment $\Lang_0$ embeds directly into this language;
all nine rules are mechanized in Lean over an idealised heap model and
exercised through an experimental prototype mode. On Pretext's four core
text-layout modules, all 105/105 units verify under the prototype with
zero unproved assumptions and all 56 guarded walker and stepper reads
restored, deriving patterns W1--W5 end-to-end while also surfacing two
upstream bounds bugs and two public-API cursor boundary checks (W6
remains an open-world API boundary assumption).
\end{abstract}

\section*{Epistemic Status}

This paper serves both as a semantic specification and as a design guide
for extending TScert's effect system. Claims fall into four distinct
tiers:
\begin{itemize}[leftmargin=1.5em,itemsep=2pt]
\item \textbf{Verified in TScert's calculus ($\Lang_0$):} Section~\ref{sec:l0}
  embeds TScert's existing footprint calculus into this language; verified
  rules cite their Lean theorem names inline.
\item \textbf{Mechanized in an idealised heap model:} The nine
  effect-variable, confined-state, scoped-callback, and region/monotonicity rules
  (\ruleName{Bound}, \ruleName{Anch}, \ruleName{Meet}, \ruleName{Inst},
  \ruleName{Confined}, \ruleName{Callback-Inv}, \ruleName{Fields-Disj},
  \ruleName{Plain-Cells-Disj}, \ruleName{Elem-Len-Extend})
  are specified over an abstract semantics in
  \lean{Verification.Paper.EffectLanguage.Rules} and proved sound over a
  heap-extension model with confined module regions and role-typed fields in
  \lean{Verification.Paper.EffectLanguage.Model} (\lean{model\_rules}).
  Transitivity of anchored effects (Proposition~\ref{prop:anchored-trans})
  follows directly from the anchoring definition.
\item \textbf{Conjectured for TScert's full calculus:}
  Conjectures~\ref{conj:regions}--\ref{conj:generic} state the region
  axioms, syntactic anchoring, and soundness of the full rule system
  against TScert's operational semantics.
\item \textbf{Experimental prototype:} An axiomatized prototype mode
  (\code{paper-rules-prototype}) applies these rules on top of TScert's
  ordinary pass under one named axiom outside the verified library,
  verifying all 105/105 units across Pretext's four core modules with
  zero unproved assumptions. Section~\ref{sec:status} states the
  prototype's soundness boundary and conservative approximations.
\end{itemize}

\section{Motivation}\label{sec:problem}

Consider the higher-order array walker of pattern P22 of \emph{Tolerant
frames in practice}~\cite{patterns}, shown below with four callers in
TScert's monomorphic footprint syntax ($\Lang_0$):

\begin{lstlisting}[language=TypeScript]
type Visitor = (i: number) => Modifies<{ values: "any"; spines: "grow" }, void>;

export function walk(xs: number[], visit: Visitor): number {
  const c = { i: 0 };
  let sum = 0;
  /** @invariant safeInteger(c.i) && 0 <= c.i */
  while (c.i < xs.length) {
    visit(c.i);
    sum += xs[c.i]!;
    c.i++;
  }
  return sum;
}

export function sumAndTrace(xs: number[]): number {
  const visited: number[] = [];
  return walk(xs, i => { visited.push(i); });
}

export function sumAndCount(xs: number[]): number {
  const seen = { n: 0 };
  return walk(xs, i => { seen.n++; });
}

export function sumAndAlias(xs: number[]): number {
  const visited = xs;
  return walk(xs, i => { visited.push(i); });
}
\end{lstlisting}

\noindent
In $\Lang_0$, TScert refuses \code{walk} at the indexed read
\code{xs[c.i]!}. The read requires $0\le c.i<\len{xs}$, where
$c.i<\len{xs}$ is established by the loop guard and is preserved by any
step within the support effect
$\Psi=\{c.i:\keep,\ xs.\kw{spine}:\grow;\ \kw{else}:\any\}$ (meaning
$c.i$ is unchanged, $xs$'s spine may only grow, and all other locations
may change arbitrarily). At \code{visit(c.i)}, however, the callback type
\code{Visitor} permits $\{\val:\any,\ \st:\grow\}$. Because $\val:\any$
covers $c.i$, the call's effect is not included in $\Psi$, so
$c.i<\len{xs}$ is dropped.

Yet \code{visit} receives only a number and cannot reach $c$, which
\code{walk} allocates after the closure already exists. The failure is an
expressiveness limit of monomorphic effects: \code{Visitor} must commit
to one fixed effect covering every caller's closure, without naming any
caller's captured variables. Because \code{sumAndCount} mutates its
captured field \code{seen.n}, the tightest closed region available in
$\Lang_0$ is the broad class $\val$, which collides with \code{walk}'s
local field $c.i$.

Section~\ref{sec:language} solves this by introducing bounded effect
variables, capture-rooted regions, and declared reach. \code{walk}
quantifies over the visitor's effect $\varepsilon$ subject to a bound
preserving $xs$'s spine; callers instantiate $\varepsilon$ with their
closure's capture-rooted effect; and \ruleName{Anch} guarantees that
$\varepsilon$ cannot modify storage allocated after the call to
\code{walk}. Section~\ref{sec:walk} works out all four callers in detail.
Scaling from a toy walker to a full text-layout library (Pretext,
Section~\ref{sec:walkers}) brings five further framing obligations:
lazy module-private caches (\ruleName{Confined}), caller invariants
needed inside non-escaping closures (\ruleName{Callback-Inv}),
object-field helper writes beside array invariants
(\ruleName{Fields-Disj}), writes to element cells of type-disjoint plain
arrays (\ruleName{Plain-Cells-Disj}), and quantified lower bounds on
nested array lengths under extension (\ruleName{Elem-Len-Extend}).

\section{The Effect Language}\label{sec:language}

The underlying heap and transition model is that of \emph{Tolerant
frames}~\cite{tf} (Sections ``What exists'' and ``Model''), strengthened
with one heap well-formedness invariant: \emph{plain-array cell
ownership} (\lean{PlainElemCellsDisjoint}), under which distinct
plain-array spines list disjoint element cells (preserved by allocation,
\code{push}, and element writes; typed-array view spines are excluded
because \code{.subarray()} views may share element cells).

\subsection{Syntax and Semantic Domains}

Commands $s$, assertions $F,P,Q,R$, and state evidence $K\in\mathcal{K}$
are imported from \emph{Tolerant frames}~\cite{tf}. Let $x,y,p\in\Var$
range over program variables ($p$ for parameters), $f,g\in\mathit{Key}$
over property keys, and $\varepsilon\in\EVar$ over effect variables.

\begin{definition}[Syntax]\label{def:syntax}
The grammar of changes, regions, effects, bounds, and polymorphic
signatures is:
{\small\[
\begin{array}{@{}l@{\quad}r@{\;\;}c@{\;\;}l@{}}
\text{Changes} & \delta,\gamma &::=& \keep\mid\init\mid\scalar\mid\extend\mid\grow\mid\any\\
\text{Field paths} & \pi &::=& x\mid\pi.f\\
\text{Regions} & r &::=& \val\mid\st\mid\vw\mid\kw{fields}\mid\pi.\kw{spine}\mid\pi.\kw{cells}\mid\pi.\kw{elemSpines}\\
 & & \mid & \pi.f\mid\pi.\kw{fields}\mid\reach{\pi}{T}\mid\modstate{M}\mid\kw{all}\mid\kw{old}\\
\text{Basic effects} & \Phi,\Psi &::=& \{r_1:\delta_1,\dots,r_k:\delta_k;\ \kw{else}:\delta_0\}\mid N(\pi_1,\dots,\pi_n)\\
\text{Effect expressions} & E &::=& \Phi\mid\varepsilon\mid E\sqcup E\mid E\sqcap E\\
\text{Bounds \& contexts} & \kappa &::=& \varepsilon\sqsubseteq\Phi \qquad\qquad C\ ::=\ \emptyset\mid C,\kappa\\
\text{Types} & T &::=& \cdots\mid(p_1:T_1,\dots,p_m:T_m)\to T\ [E\mid E^e]\\
\text{Signatures} & \Sigma &::=& \forall\varepsilon_1\dots\varepsilon_j.\ C\Rightarrow(\vec p:\vec T)\to T\ [E\mid E^e]\\
\text{Syntactic inst.} & \hat\theta &\in& \EVar\rightharpoonup\EExp
\end{array}
\]}
A named effect definition has the form $N(x_1,\dots,x_n)\triangleq\Phi$.
We write $\Delta$, $\Reg$, $\Eff$, and $\EExp$ for the sets of changes,
regions, basic effects, and effect expressions.
\end{definition}

\begin{definition}[Semantic domains]\label{def:domains}
Heaps (with contiguous domain $\{0,\dots,\len h-1\}$), environments,
closures, and semantic instantiations are:
{\small\[
\begin{array}{@{}l@{\quad}r@{\;\;}c@{\;\;}l@{}}
\text{Locations} & l &\in& \Loc = \mathbb{N}\\
\text{Values} & v &\in& \Val\\
\text{Heaps} & h &\in& \mathit{Heap}=\Loc\rightharpoonup\Val\uplus\mathit{Spine}\\
\text{Environments} & \rho,\eta &\in& \mathit{Env}=\Var\rightharpoonup\Val\\
\text{Closures} & \mathit{cl} &::=& \langle q,\eta\rangle \quad\text{(code }q\text{ with captured environment }\eta\text{)}\\
\text{Semantic inst.} & \theta &\in& \EVar\rightharpoonup\SemEff
\end{array}
\]}
where $\SemEff$ is the set of finite joins of environment-paired closed
basic effects $\langle\Phi,\rho_{\mathrm{def}}\rangle$ with
$\kw{else}:\keep$ (Definition~\ref{def:effexp}).
\end{definition}

\paragraph{Conventions.}
\begin{itemize}[leftmargin=1.5em,itemsep=2pt]
\item \textbf{Paths.} A field path $\pi$ evaluates in $(\rho,h)$ to at
  most one value $\sem{\pi}(\rho,h)$: $\rho(x)$ for a variable $x$, and
  the value stored in field $f$ of $\sem{\pi}(\rho,h)$ for $\pi.f$
  (undefined on a missing read or non-object). Inside declared reach
  (Definition~\ref{def:reach}), access-path suffixes $a$ also traverse
  array elements via $[\ast]$, evaluating $\pi a$ to the set of values
  reached (missing reads contribute nothing).
\item \textbf{Change lattice.} Each $\delta\in\Delta$ denotes, at heap
  length $b$, a reflexive and transitive binary relation $\delta_b$ on
  heap contents (\emph{Footprints}~\cite{fp}, ``Changes and
  footprints''). The linear chain
  $\keep\sqsubseteq\init\sqsubseteq\scalar\sqsubseteq\extend\sqsubseteq\grow\sqsubseteq\any$
  forms a bounded lattice $(\Delta,\sqsubseteq,\sqcup,\sqcap,\keep,\any)$
  ordered by relation inclusion.
\item \textbf{Roots.} The root of a path-based region is its leading
  variable, which may be a parameter, a local variable, or a free
  variable captured by a closure (evaluated in the closure's environment
  $\eta$).
\item \textbf{Defaults.} Omitting $\kw{else}$ defaults to
  $\delta_0=\keep$. The top effect is $\top=\{\kw{else}:\any\}$, and the
  empty (pure) effect is $\{\}=\{\kw{else}:\keep\}$. A function type
  written with a single effect $E$ defaults its exceptional effect to
  $E^e=\top$.
\end{itemize}

\subsection{Regions}

\begin{definition}[Region semantics]\label{def:regions}
Relative to an environment $\rho$, a heap $h$, and an entry heap mark
$m\le\len h$, each region $r\in\Reg$ denotes a set of locations
$\sem{r}_m(\rho,h)\subseteq\{0,\dots,\len h-1\}$ (written
$\sem{r}(\rho,h)$ when $m$ is clear from context):
\begin{itemize}[leftmargin=1.5em,itemsep=2pt]
\item $\val$, $\st$, $\vw$, and $\kw{fields}$ denote all value locations,
  all array or typed-array spine locations, all typed-array view spine
  locations, and all non-spine object property locations ($\LF$ of
  \emph{Tolerant frames}~\cite{tf}) in $h$, respectively.
\item $\pi.\kw{spine}$ denotes the spine location of the array or typed
  array (if any) at $\sem{\pi}(\rho,h)$; $\pi.\kw{cells}$ denotes the
  element locations listed by that spine in $h$; and $\pi.\kw{elemSpines}$
  denotes the plain-array spine locations of all present element values
  stored in the cells listed by $\pi$'s spine in $h$.
\item $\pi.f$ denotes the field-$f$ location of the object (if any) at
  $\sem{\pi}(\rho,h)$; $\pi.\kw{fields}$ denotes all non-spine field
  locations of that object.
\item $\reach{\pi}{T}$ denotes the declared reach of
  Definition~\ref{def:reach}.
\item $\modstate{M}$ denotes the \emph{module-private confined state} of
  a module $M$: the storage of $M$'s non-exported top-level bindings and
  their confined interior containers. A confined module region satisfies
  two module-local invariants:
  \begin{enumerate}[leftmargin=1.5em,itemsep=1pt]
  \item \emph{Boundary closure:} no reference stored outside
    $\sem{\modstate{M}}(\rho,h)$ points into $\sem{\modstate{M}}(\rho,h)$
    (\lean{Confinement.ClosedOut}), so no external root can reach
    $\modstate{M}$; and
  \item \emph{Fresh-only growth:} across any heap extension $h\to h'$,
    $\sem{\modstate{M}}(\rho,h')\subseteq\sem{\modstate{M}}(\rho,h)\cup[\len h,\len{h'})$,
    so $\modstate{M}$ never absorbs a pre-existing external location.
  \end{enumerate}
\item $\kw{all}=\{0,\dots,\len h-1\}$ denotes all allocated locations;
  and $\kw{old}=\{l\mid l<m\}$ denotes the locations already allocated at
  the entry of the enclosing invocation.
\end{itemize}
\end{definition}

\begin{example}
Executing \code{a.push(40)} updates $a$'s spine from
$\langle 4,\circ,5\rangle$ to $\langle 4,\circ,5,10\rangle$ (with
$\circ$ a hole and $10$ a fresh element location). The region
$a.\kw{spine}$ denotes the same spine location in both heaps.
\end{example}

\subsection{Declared Reach}\label{sec:reach}

Walkers over structured handles need to preserve an entire data structure
across a callback without enumerating each nested array or field. The
\emph{declared reach} $\reach{\pi}{T}$ captures all locations reachable
from $\pi$ along the keys and array elements admitted by type $T$.

\begin{definition}[Declared reach]\label{def:reach}
The set $\mathit{Acc}(T)$ of finite access-path suffixes admitted by a
type $T$ is the least set satisfying:
\begin{align*}
\mathit{Acc}(T) \;=\; \{\epsilon\}
  &\;\cup\; \{\,.f\,a \mid T\text{ object/interface declaring }f,\; a\in\mathit{Acc}(T_f)\,\}\\
  &\;\cup\; \{\,[\ast]\,a \mid T\in\{U[\,],\,\kw{readonly}\;U[\,]\},\; a\in\mathit{Acc}(U)\,\},
\end{align*}
where $T_f$ is the declared type of field $f$, union and intersection
types contribute the union of their components' access paths, and typed
arrays (\code{Uint8Array}, \code{Int32Array}, \dots) contribute only
$\{\epsilon\}$ (their interface properties are not traversed, while they
qualify as collection types below). Writing $T_a$ for the type reached
along $a\in\mathit{Acc}(T)$, the declared reach of $\pi$ at $T$ is
\begin{align*}
\reach{\pi}{T}\;=\;\bigcup_{a\in\mathit{Acc}(T)}\Bigl(
  &\{\,\pi a.f\mid T_a\text{ object/interface declaring }f\,\}\\
  &\cup\;\{\,\pi a.\kw{spine},\;\pi a.\kw{cells}\mid T_a\text{ array or typed array}\,\}\Bigr),
\end{align*}
where $T$ is the declared type of $\pi$ (or a refined type established
for $\pi$). When $a[\ast]\in\mathit{Acc}(T)$ reaches a plain array type
$T_{a[\ast]}$, the inner spine locations $\pi a.\kw{elemSpines}$ are
covered by $\pi a[\ast].\kw{spine}$.
\end{definition}

\begin{example}
Consider an abridged Pretext line-break handle:
\begin{lstlisting}[language=TypeScript]
type PreparedLineData = {
  widths: number[];
  segmentFlags: Uint8Array;
  breakableFitAdvances: (number[] | null)[];
};
\end{lstlisting}
The region $\reach{\mathit{prepared}}{\mathit{PreparedLineData}}$
contains: (i) the three field locations \code{widths},
\code{segmentFlags}, and \code{breakableFitAdvances} on \code{prepared};
(ii) the spine and element cells of \code{prepared.widths}; (iii) the
view spine and element cells of \code{prepared.segmentFlags}; (iv) the
spine and element cells of \code{prepared.breakableFitAdvances}; and (v)
the spine and element cells of every non-null inner array in
\code{prepared.breakableFitAdvances[*]} (covering
$\mathit{prepared}.\mathit{breakableFitAdvances}.\kw{elemSpines}$).
Any extra runtime property not declared by \code{PreparedLineData}, such
as \code{wholeWidth}, lies outside this reach.
\end{example}

\noindent
Because every field location's role records its property key, two field
regions with distinct keys are always disjoint regardless of aliasing.
Consequently, any field whose key is absent from the types traversed by
$\mathit{Acc}(T)$ is automatically disjoint from $\reach{\pi}{T}$, even
when stored on an object reachable from $\pi$ (used by pattern W3 in
Section~\ref{sec:walkers}).

\subsection{Region Decisions and Axioms}

From state evidence $K$, the checker decides three relations over
regions: \emph{inclusion} ($r\subseteq_K r'$), \emph{disjointness}
($r\sep_K r'$, including key-based disjointness), and \emph{stability}
($\kw{stable}_K(r)$). A region is stable when it denotes the same set of
old locations in every heap reachable by an admissible run from $K$.
Class regions are stable, as are spines and fields along a \emph{rigid
path} (a variable followed only by immutable/singleton-contract fields),
since by property (T3) of \emph{Tolerant frames} an object's keys and
field/spine locations never move once allocated.

The proof rules rely on four semantic properties of regions and their
decision procedures:
\begin{description}[leftmargin=1.5em,itemsep=1pt]
\item[(E1) Locality and bounds:] $\sem{r}_m(\rho,h)$ depends on $\rho$
  only through the root variable of $r$ (rootless regions
  $\val,\st,\vw,\kw{fields},\modstate{M},\kw{all},\kw{old}$ are
  independent of $\rho$) and contains only locations below $\len h$.
\item[(E2) Substitution:] Substituting a path $\pi$ for a root $x$
  evaluates the region at $\sem{\pi}(\rho,h)$.
\item[(E3) Sound decisions:] $r\subseteq_K r'$ and $r\sep_K r'$ imply
  set inclusion and disjointness of $\sem{r}_m(\rho,h)$ and
  $\sem{r'}_m(\rho,h)$ at every state satisfying $K$.
\item[(E4) Stability:] If $\kw{stable}_K(r)$, then for every heap $h'$
  reached by an admissible run from a state $(\rho,h)$ satisfying $K$,
  $\sem{r}_m(\rho,h')\cap\{0,\dots,\len h-1\}=\sem{r}_m(\rho,h)$.
\end{description}

\begin{conjecture}[Region axioms]\label{conj:regions}
The region semantics of Definition~\ref{def:regions} and the decision
procedures $\subseteq_K$, $\sep_K$, and $\kw{stable}_K$ satisfy
(E1)--(E4) in TScert's calculus.
\end{conjecture}

\subsection{Effects as Relations}

\begin{definition}[Effect semantics]\label{def:effect}
Given $(\rho,h)$ and mark $m\le\len h$, a basic effect
$\Phi=\{r_1:\delta_1,\ldots,r_k:\delta_k;\ \kw{else}:\delta_0\}$ assigns
each existing location $l<\len h$ a pointwise change bound in $\Delta$:
\[
\Phi_{\rho,h}(l)\;=\;
\begin{cases}
  \bigsqcup\{\,\delta_i \mid l\in\sem{r_i}_m(\rho,h)\,\} & \text{if some }r_i\text{ covers }l,\\
  \delta_0 & \text{otherwise.}
\end{cases}
\]
Compound effect expressions combine pointwise in the lattice $\Delta$:
\[
(E_1\sqcup E_2)_{\rho,h}(l) = (E_1)_{\rho,h}(l)\sqcup(E_2)_{\rho,h}(l),
\qquad
(E_1\sqcap E_2)_{\rho,h}(l) = (E_1)_{\rho,h}(l)\sqcap(E_2)_{\rho,h}(l).
\]
An effect $E$ denotes a binary relation on heaps: $h\tr{E}_\rho h'$ holds
iff $h\tr{}_\rho h'$ is an admissible transition (\emph{Tolerant
frames}) and every pre-existing location changes within its bound:
\[
h\tr{E}_\rho h' \;\iff\; h\tr{}_\rho h' \;\wedge\;
  \forall l<\len h.\;\bigl(E_{\rho,h}(l)\bigr)_{\len h}\bigl(h(l),h'(l)\bigr).
\]
Locations allocated during the step ($l\ge\len h$) are unconstrained.
\end{definition}

\noindent
Note the lattice distinction: pointwise meet $E_1\sqcap E_2$ coincides
exactly with intersection of heap-transition relations
($h\tr{E_1\sqcap E_2}_\rho h'\iff h\tr{E_1}_\rho h'\wedge h\tr{E_2}_\rho h'$),
whereas pointwise join $E_1\sqcup E_2$ is strictly larger than relational
union and allows a single composite step (such as $s_1;s_2$) to perform
$E_1$-changes on some locations and $E_2$-changes on others.

\begin{example}
$\code{a.push(40)}:\{a.\kw{spine}:\extend\}$ permits only $a$'s spine to
extend while every other pre-existing location stays $\keep$. Conversely,
$\{a.\kw{spine}:\keep;\ \kw{else}:\any\}$ constrains $a$'s spine to stay
unchanged while permitting arbitrary mutations everywhere else.
\end{example}

\begin{definition}[Subeffecting / Inclusion]\label{def:order}
Write $E\sqsubseteq_K E'$ when, for every state $(\rho,h)$ satisfying $K$
and every heap $h'$, $h\tr{E}_\rho h'$ implies $h\tr{E'}_\rho h'$. A
sound syntactic decision procedure for basic effects is given as follows:
say that $\Phi$ is \emph{within $\delta$ on $r$ under $K$} if every entry
$(r':\gamma)$ of $\Phi$ with $\neg(r\sep_K r')$ satisfies
$\gamma\sqsubseteq\delta$, and either $\delta^\Phi_0\sqsubseteq\delta$ or
$r\subseteq_K r'$ for some entry $r'$ of $\Phi$. Then
$\Phi\sqsubseteq_K\Psi$ whenever:
\begin{enumerate}[leftmargin=1.5em,itemsep=1pt]
\item $\Phi$ is within $\delta$ on $r$ for every entry $(r:\delta)$ of
  $\Psi$; and
\item if $\delta^\Psi_0\neq\any$, then $\delta^\Phi_0\sqsubseteq\delta^\Psi_0$
  and every entry $(r':\gamma)$ of $\Phi$ not contained in any region of
  $\Psi$ satisfies $\gamma\sqsubseteq\delta^\Psi_0$.
\end{enumerate}
For meets, $E_1\sqcap E_2$ is within $\delta$ on $r$ whenever either
$E_1$ or $E_2$ is; for joins, $E_1\sqcup E_2\sqsubseteq_K\Psi$ whenever
both $E_1\sqsubseteq_K\Psi$ and $E_2\sqsubseteq_K\Psi$.
\end{definition}

\begin{definition}[Fact stability]\label{def:stable}
An assertion $F$ is \emph{stable under an effect $\Psi$} if every
transition $h\tr{\Psi}_\rho h'$ preserves $F$.
\end{definition}

\begin{example}
The bound $c.i<\len{xs}$ is stable under
$\Psi=\{c.i:\keep,\ xs.\kw{spine}:\grow;\ \kw{else}:\any\}$. The effect
$\{\val:\any,\ \st:\grow\}$ is within $\grow$ on $xs.\kw{spine}$ (since
$\val\sep\st$) but not within $\keep$ on $c.i$, so it is not included in
$\Psi$.
\end{example}

\subsection{Closures, Anchoring, and Polymorphic Signatures}

\begin{definition}[Closures and function types]\label{def:closures}
A closure $\langle q,\eta\rangle$ pairs code $q$ (with parameters
$\vec p$) and an environment $\eta$ binding its free variables. The
effect stated on $q$ may reference roots in $\vec p$ or in
$\mathrm{dom}(\eta)$; at a call with actuals $\vec v$, its regions are
evaluated in $\eta[\vec p\mapsto\vec v]$. A mutable \code{let} capture
$b$ lowers to a heap cell (\emph{Tolerant frames}) and roots the field
region $b.\mathit{value}$; a \code{const} capture is a rigid variable
path. The closure has type $(\vec p:\vec T)\to T\ [E\mid E^e]$ when its
body is bounded by $E$ on normal return and $E^e$ on throw (in $\eta$)
and returns values of type $T$.
\end{definition}

\begin{example}
In \code{sumAndTrace}, the closure \code{i => \{ visited.push(i); \}} has
effect $\{\mathit{visited}.\kw{spine}:\extend\}$, rooted at the
\code{const} capture \code{visited} bound in $\eta$ to the caller's
fresh array.
\end{example}

\begin{definition}[Semantic instances and anchoring]\label{def:effexp}
A \emph{semantic effect instance} $S\in\SemEff$ is a finite join
$\bigsqcup_{j=1}^k\langle\Phi_j,\rho_j\rangle$ of environment-paired
closed basic effects, each with $\kw{else}:\keep$, evaluated pointwise
at heap $h$ and location $l<\len h$ as
$S_h(l)=\bigsqcup_{j=1}^k(\Phi_j)_{\rho_j,h}(l)$. A \emph{semantic
instantiation} $\theta:\EVar\rightharpoonup\SemEff$ maps each effect
variable to such an instance, so forwarded joins $\varepsilon_1\sqcup\Phi_2$
across distinct caller frames denote a single semantic instance without
requiring both components to share one environment. Writing
$\langle\Phi,\rho\rangle$ as $\Phi$ in $(\rho,h)$ recovers the single-frame
case.

An instance $S\in\SemEff$ is \emph{anchored} at $h_0$ (relative to the
universe $\mathcal{C}=\bigcup_M\modstate{M}$ of confined module states),
written $\Anch(S,h_0)$ (or $\Anch(\Phi,\rho,h_0)$ for
$S=\langle\Phi,\rho\rangle$), if:
\begin{enumerate}[leftmargin=1.5em,itemsep=1pt]
\item \emph{Stepwise non-widening:} for every admissible run
  $h_0\to^* h_1\to^* h_2$ and every location $l<\len{h_1}$, the
  pointwise bound never widens:
  $S_{h_2}(l)\sqsubseteq S_{h_1}(l)$; and
\item \emph{External post-entry preservation:} for every heap $h$
  reachable from $h_0$ by an admissible run and every external
  post-entry location $l\ge\len{h_0}$ with
  $l\notin\sem{\mathcal{C}}(h)$, $S_h(l)=\keep$.
\end{enumerate}
Both rigid-path regions (by (T3), $\sem{r}(\rho,h_2)=\sem{r}(\rho,h_1)$)
and confined module regions $\modstate{M}$ (by fresh-only growth,
$l<\len{h_1}\wedge l\in\sem{\modstate{M}}(h_2)\Rightarrow l\in\sem{\modstate{M}}(h_1)$)
satisfy Clause~(1), and finite joins of anchored instances remain
anchored.

A judgment or inclusion containing free effect variables holds iff it
holds under every semantic instantiation $\theta$ that maps each
quantified variable to an instance anchored at the call satisfying the
signature's bounds $C$.
\end{definition}

\begin{example}
$\{\mathit{visited}.\kw{spine}:\extend\}$, $\{\mathit{seen}.n:\any\}$,
and $\{\modstate{M}:\any\}$ are anchored at the call to \code{walk}
because the first two have \code{const}-capture roots and the third grows
only by fresh locations inside $\mathcal{C}$. By contrast, $\{\val:\any\}$
is not anchored (newly allocated external values expand $\val$), nor is
$\{\mathit{visited}.\kw{cells}:\any\}$ (pushing appends new external
element cells).
\end{example}

\noindent
All-pairs stepwise non-widening makes transitivity of anchored effects
(used by \ruleName{Loop-Eff}) an immediate consequence of
Definition~\ref{def:effexp}:

\begin{proposition}[Anchored effects are transitive]\label{prop:anchored-trans}
If $\Anch(S,h_0)$, then the transition relation $\tr{S}$ is transitive
on heaps reachable from $h_0$.
\end{proposition}

\begin{proof}
Suppose $h_0\to^* h_1\tr{S} h_2\tr{S} h_3$. By (T1) and (T2) of
\emph{Tolerant frames}~\cite{tf}, $h_1\tr{} h_3$ is admissible with
$\len{h_1}\le\len{h_2}$. For any $l<\len{h_1}$, the first step gives
$(S_{h_1}(l))_{\len{h_1}}(h_1(l),h_2(l))$ and the second gives
$(S_{h_2}(l))_{\len{h_2}}(h_2(l),h_3(l))$. By Clause~(1) of
Definition~\ref{def:effexp} on $h_0\to^* h_1\to^* h_2$,
$S_{h_2}(l)\sqsubseteq S_{h_1}(l)$, so lattice monotonicity
(\lean{Change.Holds.mono}) yields
$(S_{h_1}(l))_{\len{h_2}}(h_2(l),h_3(l))$. Transitivity of change
relations across $\len{h_1}\le\len{h_2}$ (\lean{Change.Holds.trans})
then gives $(S_{h_1}(l))_{\len{h_1}}(h_1(l),h_3(l))$, establishing
$h_1\tr{S} h_3$.
\end{proof}

\begin{conjecture}[Syntactic anchoring]\label{conj:anchored-syntax}
In TScert's calculus, any closed basic effect $\Phi$ with
$\kw{else}:\keep$ whose entries are field, fields, or spine regions along
rigid paths or confined module regions $\modstate{M}$ satisfies
$\Anch(\Phi,\rho,h_0)$ at every reachable state $(\rho,h_0)$; the same
holds for any finite join of such environment-paired effects.
\end{conjecture}

\begin{definition}[Bounds and polymorphic signatures]\label{def:constraint}
A bound $\varepsilon\sqsubseteq\Phi$ holds under $\theta$ at a state of
$K$ when $\theta(\varepsilon)\sqsubseteq_K\Phi$. In a polymorphic
signature
$\forall\vec\varepsilon.\ C\Rightarrow(p_1:T_1,\ldots,p_m:T_m)\to T\ [E\mid E^e]$,
the regions in $C$ are rooted in $\{p_1,\ldots,p_m\}$, and all effect
variables occurring in $C$, $E$, $E^e$, and $\vec T$ lie in
$\vec\varepsilon$. A function satisfies this signature if, for every
semantic instantiation $\theta$ mapping each $\varepsilon_i$ to an effect
anchored at the call satisfying $C$ there, and for all arguments of types
$\vec T[\theta]$, every normal return satisfies $E[\theta]$ and every
exceptional exit satisfies $E^e[\theta]$.
\end{definition}

\begin{example}
The polymorphic signature of \code{walk} is
\begin{align*}
\code{walk}:\;\forall\varepsilon.\;&\varepsilon\sqsubseteq\{xs.\kw{spine}:\keep;\ \kw{else}:\any\}\\
  &\Rightarrow(xs:\code{number[]},\ \mathit{visit}:(i:\code{number})\to\code{void}\ [\varepsilon])\to\code{number}\ [\varepsilon].
\end{align*}
For any anchored callback effect $\varepsilon$ preserving $xs$'s spine,
\code{walk} itself has effect at most $\varepsilon$.
\end{example}

\begin{definition}[Named effects]\label{def:named}
A named effect $N(x_1,\dots,x_n)\triangleq\Phi$ abstracts a basic effect
whose regions are rooted in $\{x_1,\dots,x_n\}$; $N(\vec\pi)$ expands to
$\Phi[\vec\pi/\vec x]$. When a module exports $N$ abstractly, it
exposes its operations' bounds in terms of $N$ while discharging the
underlying inclusions privately against $\Phi$.
\end{definition}

\section{Proof Rules}\label{sec:rules}

The proof system uses three judgments, indexed by state evidence $K$,
the enclosing function's effect bounds $C$, and the entry heap mark
$m=\len{h_0}$ (which defines $\sem{\kw{old}}_m=\{l\mid l<m\}$):
\begin{center}\small
\begin{tabularx}{\textwidth}{@{}lX@{}}
\toprule
Judgment & Meaning\\
\midrule
$K;C\vdash_m E\sqsubseteq E'$ & Subeffecting: every $E$-transition from a state of $K$ (under $C$ and mark $m$) is an $E'$-transition.\\
$K;C\vdash_m s:E\mid E^e$ & Effect typing: running $s$ from $K$ satisfies $E$ on normal return and $E^e$ on throw ($s:E$ abbreviates $s:E\mid\top$).\\
$K;C\vdash_m \hoare{P}{s}{Q}$ & Assertion preservation: Hoare-style tolerant framing of $P$ and $Q$ across $s$.\\
\bottomrule
\end{tabularx}
\end{center}

\noindent
For any effect $E$, write $E|_{\kw{old}}$ for its restriction to
pre-entry locations, defined pointwise by
$(E|_{\kw{old}})_{\rho,h}(l)=E_{\rho,h}(l)$ for $l<m$ and
$(E|_{\kw{old}})_{\rho,h}(l)=\any$ for $l\ge m$ (ignoring mutations to
storage allocated since $h_0$).
In the structural rules below, normal-exit effects are displayed
explicitly; exceptional exits compose by joining the exceptional effects
of all reachable subcommands.

\subsection{Core Structural and Framing Rules}

These rules lift the program rules of \emph{Tolerant frames}~\cite{tf}
into the unified effect language. In \ruleName{Seq-Eff}, $K'$ is the
evidence established after $s_1$, and $E_2\rightsquigarrow^{E_1} E_2'$
denotes \emph{rebasing}: $E_2'$, evaluated at entry to $s_1$,
upper-bounds every transition of $E_2$ evaluated after an $E_1$-step
(every stable region rebases to itself).

\begin{mathpar}
\inferrule*[left=Sub]
  {K;C\vdash_m s:E\mid E^e \\ K;C\vdash_m E\sqsubseteq\Psi \\ K;C\vdash_m E^e\sqsubseteq\Psi^e}
  {K;C\vdash_m s:\Psi\mid\Psi^e}
\and
\inferrule*[left=Frame]
  {K;C\vdash_m \hoare{P}{s}{Q} \\ K;C\vdash_m s:E \\ K;C\vdash_m E\sqsubseteq\Psi \\ R\text{ stable under }\Psi}
  {K;C\vdash_m \hoare{P\wedge R}{s}{Q\wedge R}}
\\
\inferrule*[left=Seq-Eff]
  {K;C\vdash_m s_1:E_1 \\ K';C\vdash_m s_2:E_2 \\ E_2\rightsquigarrow^{E_1} E_2' \\ K;C\vdash_m E_1\sqcup E_2'\sqsubseteq E}
  {K;C\vdash_m s_1;s_2:E}
\and
\inferrule*[left=If-Eff]
  {K;C\vdash_m s_1:E \\ K;C\vdash_m s_2:E}
  {K;C\vdash_m \code{if (t) s1 else s2}:E}
\and
\inferrule*[left=Loop-Eff]
  {K;C\vdash_m s:E \\ E\text{ transitive on reachable heaps}}
  {K;C\vdash_m \code{loop}\ s:E}
\\
\inferrule*[left=Call]
  {\mathit{fn}:(\vec p:\vec T)\to T\ [E\mid E^e] \\ \vec x:\vec T[\vec x/\vec p]}
  {K;C\vdash_m y=\mathit{fn}(\vec x):E[\vec x/\vec p]\mid E^e[\vec x/\vec p]}
\and
\inferrule*[left=Call-Var]
  {z:(\vec p:\vec T)\to T\ [E\mid E^e] \\ \vec x:\vec T[\vec x/\vec p]}
  {K;C\vdash_m y=z(\vec x):E[\vec x/\vec p]\mid E^e[\vec x/\vec p]}
\\
\inferrule*[left=Fun]
  {\emptyset;C\vdash_m s:E|_{\kw{old}}\mid E^e|_{\kw{old}} \\ s\text{ body of }\mathit{fn}}
  {\mathit{fn}:(\vec p:\vec T)\to T\ [E\mid E^e]}
\and
\inferrule*[left=Fresh-Eff]
  {K;C\vdash_m s:E\sqcup\{r:\delta\}\mid E^e \\ r\sep_K\kw{old}}
  {K;C\vdash_m s:E|_{\kw{old}}\mid E^e|_{\kw{old}}}
\end{mathpar}

\noindent
\status{Verified for $\Lang_0$ in Lean: \ruleName{Frame}
(\lean{respect\_sound}, \lean{frame\_step}, \lean{frame\_box});
\ruleName{Seq-Eff} (\lean{Post.compose\_rebases},
\lean{CommandSummary.seq\_sound}); \ruleName{If-Eff}
(\lean{CommandSummary.choice\_sound}); \ruleName{Loop-Eff}
(\lean{LoopEvidence.loop\_sound}); \ruleName{Call}
(\lean{call\_exits\_exact}, \lean{call\_rule}); \ruleName{Fun}
(\lean{codeObligations\_sound}, \lean{PassWalkSound}).
\ruleName{Sub}, \ruleName{Call-Var}, and \ruleName{Fresh-Eff} are not
stated as standalone Lean theorems in $\Lang_0$.}

\subsection{Subeffecting, Effect-Variable, Confined-State, and Framing Rules}

The nine rules below govern effect polymorphism, module-private confined
state, scoped callback invariants, role-typed field disjointness,
plain-array cell disjointness, and element-length monotonicity under
$\extend$, corresponding one-to-one to
\lean{BoundRule}, \lean{AnchRule}, \lean{MeetRule}, \lean{InstRule}
(mechanizing the parameter-renaming obligation \lean{RenOk}, as the
idealised model does not model closures or signatures),
\lean{ConfinedRule}, \lean{CallbackInvRule}, \lean{FieldsDisjRule},
\lean{PlainCellsDisjRule}, and \lean{ElemLenExtendRule} in
\lean{Verification.Paper.EffectLanguage.Rules}.

We use four abbreviations in the rules:
\begin{itemize}[leftmargin=1.5em,itemsep=1pt]
\item $\kw{same}_K(\Phi)$ holds when $K$ entails that $\Phi$'s root
  variables are unassigned parameters and $\Phi$ assigns every pre-entry
  location $l<m$ the same change bound as at call entry $h_0$
  (semantically, $\Same(\Phi,h_a,h_b)$ means
  $\forall l<\len{h_a}.\;\Phi_{\rho,h_b}(l)=\Phi_{\rho,h_a}(l)$).
\item $\kw{old}_{\mathcal{C}}=\kw{old}\cup\mathcal{C}$ denotes pre-entry
  locations together with confined module state (so external post-entry
  allocations lie outside $\kw{old}_{\mathcal{C}}$).
\item A region is \emph{externally rooted under $K$} if it is a concrete
  path region or declared reach whose root variable holds no reference
  into $\modstate{M}$ under $K$ (excluding rootless regions
  $\val,\st,\vw,\kw{fields},\kw{all},\modstate{M}$)
  \status{Lean: \lean{Confinement.NoPriv}}.
\item $h_0\tr{\Phi_{\mathrm{own}}\cup\Phi_{\mathrm{cb}}}^*_\rho h$
  denotes a finite run from $h_0$ to $h$ whose steps each satisfy
  $h_a\tr{\Phi_{\mathrm{own}}}_\rho h_b$ or
  $h_a\tr{\Phi_{\mathrm{cb}}}_\rho h_b$; and $\mathrm{prefix}(\pi)$
  denotes the field locations traversed along path $\pi$.
\end{itemize}

\begin{mathpar}
\inferrule*[left=Bound]
  {(\varepsilon\sqsubseteq\Phi)\in C \\ \kw{same}_K(\Phi)}
  {K;C\vdash_m \varepsilon\sqsubseteq\Phi}
\and
\inferrule*[left=Anch]
  {\varepsilon\text{ bound by enclosing signature}}
  {K;C\vdash_m \varepsilon\sqsubseteq\{\kw{old}_{\mathcal{C}}:\any\}}
\and
\inferrule*[left=Meet]
  {K;C\vdash_m E\sqsubseteq E_1 \\ K;C\vdash_m E\sqsubseteq E_2}
  {K;C\vdash_m E\sqsubseteq E_1\sqcap E_2}
\\
\inferrule*[left=Inst]
  {\mathit{fn}:\forall\vec\varepsilon.\;C'\Rightarrow(\vec p:\vec T)\to T\ [E\mid E^e] \\
   \forall\varepsilon\in\vec\varepsilon.\;\hat\theta(\varepsilon)\text{ anchored under }C \\
   K;C\vdash_m C'[\vec x/\vec p][\hat\theta] \\
   \vec x:\vec T[\vec x/\vec p][\hat\theta]}
  {K;C\vdash_m y=\mathit{fn}(\vec x):E[\vec x/\vec p][\hat\theta]\mid E^e[\vec x/\vec p][\hat\theta]}
\\
\inferrule*[left=Confined]
  {\modstate{M}\text{ confined in }M \\
   \Phi\text{ has }\kw{else}:\any\text{ and every entry of }\Phi\text{ is externally rooted under }K}
  {K;C\vdash_m \{\modstate{M}:\any\}\text{ anchored}\qquad K;C\vdash_m \{\modstate{M}:\any\}\sqsubseteq\Phi}
\\
\inferrule*[left=Callback-Inv]
  {P(h_0) \quad P\text{ stable under }\Psi \\
   \forall h_a,h_b.\;h_a\tr{\Psi}_\rho h_b\Rightarrow\Same(\Psi,h_a,h_b) \\
   K;C\vdash_m \Phi_{\mathrm{own}}\sqsubseteq\Psi \quad \Phi_{\mathrm{own}}\text{ anchored} \\
   K;C\vdash_m \Phi_{\mathrm{cb}}\sqsubseteq\Psi \quad \Phi_{\mathrm{cb}}\text{ anchored} \quad \text{callback does not escape}}
  {\forall h.\;h_0\tr{\Phi_{\mathrm{own}}\cup\Phi_{\mathrm{cb}}}^*_\rho h\Rightarrow P(h)}
\\
\inferrule*[left=Fields-Disj]
  {\Phi\text{ has }\kw{else}:\any \\
   \text{every constrained entry of }\Phi\text{ is non-field (spines, element cells, or element spines)}}
  {K;C\vdash_m \{\kw{fields}:\any;\ \kw{else}:\keep\}\sqsubseteq\Phi}
\\
\inferrule*[left=Plain-Cells-Disj]
  {\pi\text{ plain array} \quad \Phi\text{ has }\kw{else}:\any \\
   \text{every constrained entry of }\Phi\text{ has the form }(\pi'.\kw{cells}:\gamma)\\
   \text{with }\pi'\text{ plain array and }\pi.\kw{spine}\sep_K\pi'.\kw{spine}}
  {K;C\vdash_m \{\pi.\kw{cells}:\delta;\ \kw{else}:\keep\}\sqsubseteq\Phi}
\\
\inferrule*[left=Elem-Len-Extend]
  {P:\mathbb{N}\to\mathrm{Prop}\text{ monotone} \\
   \begin{aligned}
     K;C\vdash_m \Phi\sqsubseteq\{&\mathrm{prefix}(\pi):\keep,\ \pi.\kw{spine}:\keep,\cr
       &\pi.\kw{cells}:\keep,\ \pi.\kw{elemSpines}:\extend;\ \kw{else}:\any\}
   \end{aligned}}
  {(\forall i<\len\pi.\;\pi[i]\text{ plain array}\Rightarrow P(\len{\pi[i]}))\text{ is stable under }\Phi}
\end{mathpar}

\begin{description}[leftmargin=1.5em,itemsep=2pt]
\item[\ruleName{Bound}] An effect variable is bounded by its declared
  constraint $\Phi$ at any later state where $\kw{same}_K(\Phi)$ holds.
  A sufficient condition is that $\Phi$'s roots are unassigned parameters,
  $\Phi$ traverses only regions that $\Phi$ itself keeps unchanged, and
  every step since $h_0$ lies within $\Phi$.
\item[\ruleName{Anch}] Every semantic instance of $\varepsilon$ is
  anchored at $h_0$, so by Clause~(2) of Definition~\ref{def:effexp} it
  leaves every external post-entry location ($l\ge m$,
  $l\notin\mathcal{C}$) unchanged; equivalently, it is included in
  $\{\kw{old}_{\mathcal{C}}:\any;\ \kw{else}:\keep\}$.
\item[\ruleName{Meet}] Inclusion distributes over meets, allowing a
  callback call to combine \ruleName{Bound} (preserving the handle) with
  \ruleName{Anch} (preserving freshly allocated locals) and chain
  additional constraints.
\item[\ruleName{Inst}] Instantiates a polymorphic signature at actual
  arguments $\vec x$ (substituted before $\hat\theta$, whose images live
  in the caller's scope). Four common syntactic instantiations arise:
  \begin{enumerate}[leftmargin=1.5em,itemsep=1pt]
  \item \emph{Concrete anchored effect:} $\hat\theta(\varepsilon)=\Phi$
    with $\kw{else}:\keep$ (including $\{\modstate{M}:\any\}$ or its join
    with capture-rooted entries), checked against $C'[\vec x/\vec p]$ by
    Definition~\ref{def:order} and \ruleName{Confined}.
  \item \emph{Omitted optional callback:} if every parameter mentioning
    $\varepsilon$ receives \code{undefined}, $\hat\theta(\varepsilon)=\{\}$
    trivially satisfies every bound.
  \item \emph{Polymorphic forwarding:} inside an enclosing function
    quantifying $\varepsilon$, setting $\hat\theta(\varepsilon')=\varepsilon$
    is anchored under $C$, and its bound check reduces to showing the
    caller's bound on $\varepsilon$ (read at $\vec x$), combined with
    \ruleName{Anch} and \ruleName{Meet} on own locations of caller-allocated
    fresh objects, is included in the callee's bound on $\varepsilon'$.
  \item \emph{Forwarded join:} setting
    $\hat\theta(\varepsilon')=\varepsilon\sqcup\Phi$ (for instance
    $\varepsilon\sqcup\{\modstate{M}:\any\}$ when an adapter invokes a
    caching helper before forwarding to its caller's callback) is anchored
    whenever $\Phi$ is anchored, and discharges against $C'[\vec x/\vec p]$
    componentwise.
  \end{enumerate}
  By \ruleName{Sub}, a caller may also over-approximate an instantiated
  effect by substituting its bound in place of $\hat\theta(\varepsilon)$.
\item[\ruleName{Confined}] A module-private confined state region
  $\modstate{M}$ is anchored at every state (fresh-only growth keeps
  external post-entry allocations unchanged) and is disjoint from every
  externally rooted region---both concrete path regions and declared
  reaches $\reach{x}{T}$---by boundary closure. Consequently,
  $\{\modstate{M}:\any\}$ automatically satisfies any externally rooted
  callback bound with $\kw{else}:\any$, and preserves external handle
  invariants across helper calls
  \status{Lean: \lean{Confinement.ClosedOut}, \lean{closedOut\_region\_disjoint}, \lean{closedOut\_reach\_disjoint}}.
\item[\ruleName{Callback-Inv}] A caller invariant $P$ that holds at call
  entry $h_0$ and is preserved by a self-stabilizing support effect $\Psi$
  remains true at every heap reached by interleaving the callee's own steps
  ($\Phi_{\mathrm{own}}$) and the callback's steps ($\Phi_{\mathrm{cb}}$)
  whenever both are anchored and within $\Psi$ at $h_0$. When the callback
  does not escape the callee's dynamic extent, every invocation of the
  callback occurs in such a heap, so the callback body may assume $P$ over
  its captured variables on entry.
\item[\ruleName{Fields-Disj}] Non-spine object property locations
  ($\kw{fields}$) are disjoint by allocation role from array and view
  spines, array element cells, and element spines. Consequently, a helper
  whose effect is $\{\kw{fields}:\any;\ \kw{else}:\keep\}$ (or
  $\{\kw{fields}:\any,\ \modstate{M}:\any;\ \kw{else}:\keep\}$) preserves
  every spine, element-cell, and element-spine invariant
  \status{Lean: \lean{Role.field}, \lean{fieldClass\_region\_disjoint}}.
\item[\ruleName{Plain-Cells-Disj}] Under plain-array cell ownership,
  distinct plain-array spines own disjoint element cells; typed-array
  views are excluded because \code{.subarray()} views may share cells.
  When static types or allocation separation guarantee two plain-array
  paths have apart spines, writes to one array's element cells preserve
  facts over the other array's element cells
  \status{Lean: \lean{PlainElemCellsDisjoint}, \lean{plainArray\_cells\_disjoint}, \lean{typesPathSpinesApart\_sound}}.
\item[\ruleName{Elem-Len-Extend}] Any $\extend$-bounded transition on a
  plain-array spine leaves it an array spine of at least the same length.
  Hence an effect that keeps an outer array's path, spine, and element
  cells while bounding its element spines by $\extend$ preserves every
  monotone lower bound on element lengths (such as
  $\forall i<\len\pi.\;\pi[i]=\code{null}\vee 1\le\len{\pi[i]}$)
  \status{Lean: \lean{extend\_arr\_back}}.
\end{description}

\begin{theorem}[Idealised model soundness]\label{thm:model}
Over the idealised heap-extension model of
\lean{Verification.Paper.EffectLanguage.Model} (built on
\lean{Verification.Semantics.Footprints} and
\lean{Verification.Semantics.Confinement}), for any universe $C$ of
confined module regions (\lean{ModuleRegion}) and role-typed field class
$F$ (\lean{FieldClass}):
\begin{enumerate}[leftmargin=1.5em,itemsep=1pt]
\item Boundary closure implies disjointness of $\modstate{M}$ from every
  externally rooted concrete region (\lean{closedOut\_region\_disjoint},
  \lean{modDisjoint\_of\_external}) and from every externally rooted
  declared reach $\reach{x}{T}$ (\lean{closedOut\_reach\_disjoint}).
\item Joining confined module regions into an anchored effect preserves
  anchoring and inclusion (\lean{eff1\_add\_modules\_anchored},
  \lean{eff1\_add\_modules\_within}).
\item Non-spine object fields are disjoint from spines, element cells, and
  element spines (\lean{fieldClass\_region\_disjoint},
  \lean{fieldsDisj\_of\_nonField}); distinct plain-array spines have
  disjoint element cells (\lean{plainArray\_cells\_disjoint},
  \lean{plainCellsDisj\_of\_apart}); and $\extend$ preserves monotone
  element-length lower bounds (\lean{extend\_arr\_back}).
\item All nine rules \ruleName{Bound}, \ruleName{Anch}, \ruleName{Meet},
  \ruleName{Inst} (its parameter-renaming obligation \lean{RenOk}),
  \ruleName{Confined}, \ruleName{Callback-Inv}, \ruleName{Fields-Disj},
  \ruleName{Plain-Cells-Disj}, and \ruleName{Elem-Len-Extend} are sound
  (\lean{model\_rules}).
\end{enumerate}
\end{theorem}

\subsection{Primitive Step Rules}

For a field \code{x.f}, a plain array \code{x} (a typed-array element
write \code{x[i] = a} carries effect $\{\vw:\grow\}$ because
\code{.subarray()} views may share cells), and a fresh variable \code{y}:

\begin{mathpar}
\inferrule*[left=Field-Write]{ }{K;C\vdash_m \code{x.f = a}:\{x.f:\any\}}
\and
\inferrule*[left=Field-Write-Scalar]{x.f\text{ and }a\text{ primitive}}{K;C\vdash_m \code{x.f = a}:\{x.f:\scalar\}}
\\
\inferrule*[left=Elem-Write]{x\text{ plain array}}{K;C\vdash_m \code{x[i] = a}:\{x.\kw{spine}:\grow,\ x.\kw{cells}:\any\}}
\and
\inferrule*[left=Dense-Elem-Write]{x:\code{Dense<}T\code{[]>}\quad K\vdash 0\le i < x.\code{length}}{K;C\vdash_m \code{x[i] = a}:\{x.\kw{cells}:\any\}}
\\
\inferrule*[left=Push]{x\text{ plain array}}{K;C\vdash_m \code{x.push(a)}:\{x.\kw{spine}:\extend\}}
\\
\inferrule*[left=Alloc]{ }{K;C\vdash_m \code{y = \{f: a\}}:\{\}\qquad K;C\vdash_m \code{y = [a]}:\{\}}
\and
\inferrule*[left=Native]{\mathit{nat}:(\vec p:\vec T)\to T\ [\Phi_{\mathit{nat}}]}{K;C\vdash_m y=\mathit{nat}(\vec x):\Phi_{\mathit{nat}}[\vec x/\vec p]}
\end{mathpar}

\noindent
\status{Verified in Lean by \lean{fieldWriteFrame} and
\lean{fieldWriteStep\_sound}, \lean{ExpressionEffects.scalarWrite\_sound},
\lean{elemWriteFrame} and \lean{elemWriteStep\_sound},
\lean{ExpressionEffects.elem\_var\_write\_dense\_core} and
\lean{IndexWrites.boundedWriteFrame\_sound} (with
\lean{IndexWrites.denseWriteFootprint}),
\lean{receiverExtendFrame} (under \lean{ReceiverAppends}), and
\lean{alloc\_post}. Native contracts are audited specifications
(\emph{Footprints}~\cite{fp}).}

\begin{conjecture}[Full calculus soundness]\label{conj:generic}
The rules of Section~\ref{sec:rules} are sound for the full operational
semantics of TScert's calculus, with heap well-formedness extended by
plain-array cell ownership (\lean{PlainElemCellsDisjoint}) for
\ruleName{Plain-Cells-Disj}.
\end{conjecture}

\section{Surface Syntax in TypeScript}\label{sec:annotations}

Effects are written in TypeScript using \code{Modifies}~\cite{ann},
\code{Effects}, \code{AtMost}, and \code{Footprint}, all erased at
runtime:

\begin{center}\small
\begin{tabularx}{\textwidth}{@{}lX@{}}
\toprule
Surface syntax & Denotation\\
\midrule
\code{values: "any"}, \code{spines: "grow"} & $\val:\any$, $\st:\grow$\\
\code{views: "grow"}, \code{fields: "any"} & $\vw:\grow$, non-spine object fields $\kw{fields}:\any$\\
\code{"visited[*]": "extend"}, \code{"xs[\#]": "any"} & $\mathit{visited}.\kw{spine}:\extend$, element cells $xs.\kw{cells}:\any$\\
\code{"flow.wholeWidth": "any"}, \code{"x.*": "any"} & field $\mathit{flow}.\mathit{wholeWidth}:\any$, all fields $x.\kw{fields}:\any$\\
\code{"prepared.**": "keep"} & declared reach $\reach{\mathit{prepared}}{T}:\keep$ ($T$ declared type of \code{prepared})\\
\code{"@state": "any"} & module-private confined state $\modstate{M}:\any$ of the defining module $M$\\
\code{else: "any"} & default change $\delta_0$ (\code{"keep"} when omitted)\\
\code{onThrow: true} & same effect at exceptional exit ($E^e=E$)\\
\code{Modifies<P, T>} & return type $T$ with effect $P$ (read by ordinary TScert and the prototype)\\
\code{Effects<P, T>} & return type $T$ with effect-language footprint $P$ (supports \code{"@state"}, \code{fields}, \code{views}, \code{p.**}, joins; erased by ordinary TScert)\\
\code{E extends Footprint} & effect type parameter constraint for callback types\\
\code{<E extends AtMost<P>>} & quantified effect variable $\varepsilon$ with bound $\varepsilon\sqsubseteq P$\\
\code{Modifies<E, T>}, \code{Join<E, P>} & effect $\varepsilon$, and pointwise join $\varepsilon\sqcup P$\\
\code{f<P>(...)} & explicit syntactic instantiation $\hat\theta(\varepsilon)=P$\\
\code{type N<x extends string> = ...} & named effect $N(x)\triangleq P$\\
\bottomrule
\end{tabularx}
\end{center}

\noindent
Because TypeScript treats \code{Modifies<P, T>} as \code{T} and erases
\code{P}, type inference does not reconstruct effect arguments; call
sites write \code{f<P>(...)} explicitly (except for omitted optional
callbacks and forwarded variables).

\section{Worked Example: \code{walk} and its Callers}\label{sec:walk}

Rewriting Section~\ref{sec:problem} with a bounded effect variable:

\begin{lstlisting}[language=TypeScript]
type Visitor<E extends Footprint> = (i: number) => Modifies<E, void>;

export function walk<E extends AtMost<{ "xs[*]": "keep"; else: "any" }>>(
  xs: number[], visit: Visitor<E>,
): Modifies<E, number> {
  // same body as Section 1
}

export function sumAndTrace(xs: number[]): Modifies<{}, number> {
  const visited: number[] = [];
  const trace = (i: number): Modifies<{ "visited[*]": "extend" }, void> => {
    visited.push(i);
  };
  return walk<{ "visited[*]": "extend" }>(xs, trace);
}

export function sumAndCount(xs: number[]): Modifies<{}, number> {
  const seen = { n: 0 };
  const count = (i: number): Modifies<{ "seen.n": "any" }, void> => {
    seen.n++;
  };
  return walk<{ "seen.n": "any" }>(xs, count);
}

export function sumAndAlias(xs: number[]): Modifies<{ "xs[*]": "extend" }, number> {
  const visited = xs;
  const trace = (i: number): Modifies<{ "visited[*]": "extend" }, void> => {
    visited.push(i);
  };
  return walk<{ "visited[*]": "extend" }>(xs, trace);
}
\end{lstlisting}

\begin{description}[leftmargin=1.5em,itemsep=3pt]
\item[\code{walk} (checked once polymorphically).]
  At entry, $C=\{\varepsilon\sqsubseteq\{xs.\kw{spine}:\keep;\ \kw{else}:\any\}\}$.
  Allocating \code{const c = \{ i: 0 \}} places $c$'s fields outside
  $\kw{old}_{\mathcal{C}}$ (\ruleName{Alloc}, since $c$ is a fresh
  external local outside both $\kw{old}$ and confined module state
  $\mathcal{C}$). The loop guard yields $c.i<\len{xs}$, stable under
  $\Psi=\{c.i:\keep,\ xs.\kw{spine}:\grow;\ \kw{else}:\any\}$. At
  \code{visit(c.i)}, \ruleName{Call-Var} gives effect $\varepsilon$. By
  \ruleName{Bound} (using (T3) on parameter $xs$) and \ruleName{Anch},
  \ruleName{Meet} yields
  $\varepsilon\sqsubseteq\{xs.\kw{spine}:\keep;\ \kw{else}:\any\}\sqcap\{\kw{old}_{\mathcal{C}}:\any\}$.
  This meet is included in $\Psi$: the second component keeps $c.i$
  (disjoint from $\kw{old}_{\mathcal{C}}$), the first keeps
  $xs.\kw{spine}$ (within $\grow$), and $\Psi$'s $\kw{else}$ is $\any$.
  Hence \ruleName{Frame} preserves $c.i<\len{xs}$ and \code{xs[c.i]!} is
  in bounds. Finally, \ruleName{Seq-Eff}, \ruleName{Loop-Eff} (via
  Proposition~\ref{prop:anchored-trans}), \ruleName{Fresh-Eff} (removing
  writes to $c$), and \ruleName{Fun} establish \code{Modifies<E, number>}.
\item[\code{sumAndTrace} (accepted).]
  \ruleName{Push} and \ruleName{Fun} give \code{trace} the anchored
  effect $\{\mathit{visited}.\kw{spine}:\extend\}$. Since \code{visited}
  is freshly allocated, its spine is disjoint from $xs.\kw{spine}$, so
  \ruleName{Inst} discharges the bound. \ruleName{Fresh-Eff} drops the
  mutation to \code{visited}, yielding \code{Modifies<\{\}, number>}.
\item[\code{sumAndCount} (accepted).]
  The anchored instance $\{\mathit{seen}.n:\any\}$ touches only a value
  field (disjoint from any spine), discharging the bound by
  \ruleName{Inst}. No monomorphic effect in $\Lang_0$ can express this
  without widening to $\val:\any$.
\item[\code{sumAndAlias} (refused).]
  Here \code{visited} aliases \code{xs}, so $\mathit{visited}.\kw{spine}$
  and $xs.\kw{spine}$ coincide with $\extend\not\sqsubseteq\keep$, and
  \ruleName{Inst} refuses the call. Weakening \code{walk}'s bound to
  \code{\{ "xs[*]": "grow"; else: "any" \}} accepts \code{sumAndAlias} as
  well: $c.i<\len{xs}$ remains stable under spine growth (the loop then
  visits appended elements until \code{push} throws at maximal array
  length, which TScert permits).
\end{description}

\section{Case Study: Pretext's Line Walkers}\label{sec:walkers}

We test the language against the six callback patterns W1--W6 and the
four core text-layout modules of Pretext (commit \code{0dafecb3}). Define
the named effect preserving a line-break handle's declared reach:
\[
\mathit{Geom}(p)\;\triangleq\;\{\reach{p}{\mathit{PreparedLineData}}:\keep;\ \kw{else}:\any\},
\]
written in TypeScript as
\code{type Geom<p extends string> = Record<`\$\{p\}.**`, "keep"> \& \{ else: "any" \}}
(expanded inline in the prototype, which requires literal effect types;
Section~\ref{sec:status}). In the end-to-end Pretext suite below, all
internal helpers are verified directly: \code{stepPreparedSimpleLineGeometry}
writes only its cursor's fields and confined module state (\code{"@state"},
the lazy engine-profile cache in \code{measurement.ts}), while
\code{getInternalPrepared}, \code{normalizeMaxWidth}, and
\code{createLayoutLineRange} leave pre-existing storage unchanged.

\subsection{W1 and W2: Local Cursor (\ruleName{Anch}) and Handle Reach (\ruleName{Bound})}

\code{walkPreparedLinesRaw} allocates a local cursor, steps through the
handle's segments, updates \code{stats}, and invokes an optional visitor
on each line:

\begin{lstlisting}[language=TypeScript]
type InternalLineVisitor<E extends Footprint> = (
  width: number, startSeg: number, startGraph: number,
  endSeg: number, endGraph: number,
) => Modifies<E, void>;

export function walkPreparedLinesRaw<E extends AtMost<Geom<"prepared">>>(
  prepared: PreparedLineData, maxWidth: number, onLine?: InternalLineVisitor<E>,
  stats: LineStats = { lineCount: 0, maxLineWidth: 0 },
): Effects<Join<E, { "stats.lineCount": "any"; "stats.maxLineWidth": "any"; "@state": "any" }>, number> {
  const cursor: LayoutCursor = { segmentIndex: 0, graphemeIndex: 0 };
  const { segmentFlags } = prepared;
  const segmentCount = prepared.widths.length;
  while (true) {
    let startSegmentIndex = cursor.segmentIndex;
    while (startSegmentIndex < segmentCount &&
           consumesAtLineStart(segmentFlags[startSegmentIndex]!)) {
      startSegmentIndex++;
    }
    if (startSegmentIndex >= segmentCount) return stats.lineCount;
    const startGraphemeIndex = cursor.graphemeIndex;
    cursor.segmentIndex = startSegmentIndex;
    const width = stepPreparedSimpleLineGeometry(prepared, cursor, maxWidth);
    stats.lineCount++;
    if (width > stats.maxLineWidth) stats.maxLineWidth = width;
    onLine?.(width, startSegmentIndex, startGraphemeIndex,
             cursor.segmentIndex, cursor.graphemeIndex);
  }
}
\end{lstlisting}

\noindent
Across \code{onLine?.(...)}:
\begin{itemize}[leftmargin=1.5em,itemsep=2pt]
\item \textbf{W1 (local cursor):} \code{cursor} is allocated after
  entry, so its fields lie outside $\kw{old}_{\mathcal{C}}$ and survive by
  \ruleName{Anch}.
\item \textbf{W2 (handle reach):} The handle invariants (such as
  $\len{\mathit{prepared}.\mathit{widths}}=\len{\mathit{prepared}.\mathit{segmentFlags}}$)
  and the alias linking \code{segmentFlags} to
  \code{prepared.segmentFlags} are stable under
  $\mathit{Geom}(\mathit{prepared})$ and survive by \ruleName{Bound}.
  The walker's own writes to \code{stats} and \code{cursor} use keys
  absent from \code{PreparedLineData}, and
  \code{stepPreparedSimpleLineGeometry}'s cache writes lie in
  $\modstate{M}$, so all body steps are disjoint from the handle's reach
  even if \code{stats} aliased \code{prepared}, preserving
  $\kw{same}_K(\mathit{Geom}(\mathit{prepared}))$.
\end{itemize}

\subsection{W3: Key-Disjoint Writes on a Container}

\code{findWholeLine} walks \code{flow.data} while writing summary fields
directly onto the enclosing \code{flow} object:

\begin{lstlisting}[language=TypeScript]
function findWholeLine(flow: InternalPreparedRichInline): InternalPreparedRichInline {
  const lineCount = walkPreparedLinesRaw<{
    "flow.wholeWidth": "any"; "flow.wholeStart": "any"; "flow.wholeEnd": "any";
  }>(flow.data, Number.POSITIVE_INFINITY,
    (width, startSeg, _startGraph, endSeg) => {
      flow.wholeWidth = width;
      flow.wholeStart = startSeg;
      flow.wholeEnd = endSeg;
    });
  if (lineCount !== 1) flow.wholeWidth = null;
  return flow;
}
\end{lstlisting}

\noindent
The callback's anchored effect is
$\Phi_3=\{\mathit{flow}.\mathit{wholeWidth}:\any,\ \mathit{flow}.\mathit{wholeStart}:\any,\ \mathit{flow}.\mathit{wholeEnd}:\any\}$.
Because none of those three keys appears in the declared reach of
\code{PreparedLineData}, $\Phi_3$ is key-disjoint from
$\reach{(\mathit{flow}.\mathit{data})}{\mathit{PreparedLineData}}$, so
\ruleName{Inst} succeeds without proving that \code{flow} and
\code{flow.data} are distinct objects.

\subsection{W4: Polymorphic Forwarding and Abstract Named Effects}

The wrapper \code{walkLineRanges} quantifies over its caller's
effect $\varepsilon$ and forwards it to \code{walkPreparedLinesRaw}:

\begin{lstlisting}[language=TypeScript]
export function walkLineRanges<E extends AtMost<Layout<"prepared">>>(
  prepared: PreparedText, maxWidth: number,
  onLine: (line: LayoutLineRange) => Modifies<E, void>,
): Effects<Join<E, { "@state": "any" }>, number> {
  return walkPreparedLinesRaw<E>(
    getInternalPrepared(prepared), normalizeMaxWidth(maxWidth),
    (w, sSeg, sGraph, eSeg, eGraph): Modifies<E, void> => {
      onLine(createLayoutLineRange(w, sSeg, sGraph, eSeg, eGraph));
    });
}
\end{lstlisting}

\noindent
The adapter closure allocates a line range and invokes the captured
\code{onLine}, giving effect $\varepsilon$ by \ruleName{Seq-Eff}. At the
inner call, \ruleName{Inst} instantiates $\hat\theta(\varepsilon')=\varepsilon$:
the obligation $\varepsilon\sqsubseteq\mathit{Geom}(t)$ (for
$t=\code{getInternalPrepared(prepared)}$) discharges from
$\varepsilon\sqsubseteq\mathit{Layout}(\mathit{prepared})$ once the
module expands its private definition
$\mathit{Layout}(p)\triangleq\mathit{Geom}(p)$ and links $t$ to
\code{prepared}. Omitting \code{stats} allocates a fresh default object
whose field mutations drop by \ruleName{Fresh-Eff}, leaving the returned
effect $\varepsilon\sqcup\{\modstate{M}:\any\}$.

\subsection{W5 and W6: Confined Module Caches and Open-World Boundaries}

\begin{description}[leftmargin=1.5em,itemsep=2pt]
\item[W5 (collector with confined module caches).]
  \code{layoutWithLines} collects lines into a fresh array \code{lines}
  by calling \code{createLayoutLine} on \code{prepared} inside the
  visitor passed to \code{walkPreparedLinesRaw}:
  \[
    \hat\theta(\varepsilon)\;=\;\{\mathit{lines}.\kw{spine}:\extend,\ \modstate{M}:\any\}.
  \]
  Transitively, \code{createLayoutLine} populates module-private lazy
  caches: a \code{WeakMap} (\code{graphemeEndCaches} keyed by
  \code{PreparedText}), unpacked lookup tables (\code{graphemeRules},
  \code{breakRules}), and a memoized engine profile
  (\code{cachedEngineProfile}). Each defining module satisfies boundary
  closure and fresh-only growth locally: every cache binding is
  unexported, initialized with fresh objects or primitives, never
  returned or stored into parameters or exports, and never mutates
  values read from a parameter. Thus \code{createLayoutLine} carries
  effect \code{Effects<\{ "@state": "any" \}, LayoutLine>}. At
  \code{walkPreparedLinesRaw}, \ruleName{Inst} and \ruleName{Confined}
  discharge both components against $\mathit{Geom}(\mathit{prepared})$:
  \code{lines} is fresh and disjoint from \code{prepared}, and
  $\modstate{M}$ is disjoint from
  $\reach{\mathit{prepared}}{\mathit{PreparedLineData}}$ and anchored.
  Similarly, \code{walkRichInlineLineRanges} wraps its caller's callback
  with a caching helper and forwards the joined effect
  \code{Join<E, \{ "@state": "any" \}>} ($\varepsilon\sqcup\{\modstate{M}:\any\}$)
  to \code{walkPreparedLinesRaw}, where $\reach{(\mathit{prepared}.\mathit{data})}{\mathit{PreparedLineData}}\subseteq\reach{\mathit{prepared}}{\mathit{PreparedRichInline}}$.
\item[W6 (unanalysed callers).] Callers verified by TScert discharge
  \code{walkLineRanges}'s bound via \ruleName{Inst}; unanalysed external
  callers assume the bound at the API boundary, exactly as with ordinary
  preconditions.
\end{description}

\subsection{Additional Pretext Patterns and End-to-End Evaluation}\label{sec:eval}

Verifying Pretext's four core modules end-to-end exercises the remaining
four proof rules alongside W1--W5:
\begin{itemize}[leftmargin=1.5em,itemsep=2pt]
\item \textbf{\ruleName{Callback-Inv}:} In \code{layoutWithLines} and
  \code{walkRichInlineLineRanges}, the inline visitor passed to
  \code{walkPreparedLinesRaw} reads captured handle arrays
  (\code{prepared}, \code{flow}); because the callback does not escape
  and both the walker and visitor preserve the handle's reach, the
  caller's handle invariants hold on entry to every callback invocation.
\item \textbf{\ruleName{Fields-Disj}:} Font-measurement, Han-kerning, and
  DOM style helpers (\code{han-kerning.ts}, \code{measurement.ts},
  \code{prepare.ts}) mutate object properties (\code{fields: "any"}) while
  array preparation loops maintain length and element-cell invariants over
  output arrays.
\item \textbf{\ruleName{Plain-Cells-Disj} and \ruleName{Elem-Len-Extend}:}
  The preparation builders (\code{prepare.ts}, \code{rich-inline.ts})
  populate parallel plain arrays (\code{widths}, \code{segments}, and
  the nested array \code{breakableFitAdvances}, written $b$) while
  establishing the handle's quantified invariant
  $\forall i<\len b.\;b[i]=\code{null}\vee 1\le\len{b[i]}$;
  type-disjoint plain-array cell writes preserve sibling arrays' element
  invariants, and appending to inner fit-advance arrays under $\extend$
  preserves the non-empty lower bound.
\end{itemize}

\begin{center}\small
\begin{tabularx}{\textwidth}{@{}lrrX@{}}
\toprule
Pretext module & Units accepted & Assumptions & Key obligations discharged\\
\midrule
\code{src/line-break.ts} & 17 / 17 & 0 & All 56 guarded walker and stepper reads restored; W1, W2, \ruleName{Confined}\\
\code{src/layout.ts} & 18 / 18 & 0 & W4, W5, \ruleName{Callback-Inv}; 1 public-API cursor boundary check (\code{stepNextLine})\\
\code{src/rich-inline.ts} & 38 / 38 & 0 & W3, W5, \ruleName{Callback-Inv}, \ruleName{Plain-Cells-Disj}; 1 cursor boundary check (\code{layoutNextRichInlineLineRange})\\
\code{src/prepare.ts} & 32 / 32 & 0 & Full handle invariants, \ruleName{Fields-Disj}, \ruleName{Plain-Cells-Disj}, \ruleName{Elem-Len-Extend}\\
\midrule
\textbf{Total} & \textbf{105 / 105} & \textbf{0} & \textbf{56 guards restored; 2 upstream bounds bugs + 2 cursor boundary checks}\\
\bottomrule
\end{tabularx}
\end{center}

\noindent
Verification also surfaced two genuine upstream bounds bugs in Pretext's
supporting modules, both reproduced directly under Node:
(i) \code{getGraphemeStart} in \code{line-text.ts} reads
\code{ends[graphemeIndex - 1]!}, which returns \code{undefined} typed as
\code{number} when a line range from a longer text is materialized on a
shorter text past a segment's last grapheme (producing the whole segment
or \code{NaN} source offsets); and
(ii) \code{getCanvasFont} in \code{measurement.ts} indexes
\code{families[...]} out of bounds when a profile lists fewer generic
families than keywords, interpolating \code{"undefined"} into a CSS font
family name. Separately, the two public streaming entry points
(\code{stepNextLine} in \code{layout.ts} and
\code{layoutNextRichInlineLineRange} in \code{rich-inline.ts}) add
runtime cursor boundary checks validating caller-supplied cursors; these
are public-API boundary guards rather than upstream bugs on
Pretext-produced cursors.

\section{\texorpdfstring{Embedding TScert's Verified Fragment $\Lang_0$}{Embedding TScert's Verified Fragment L0}}\label{sec:l0}

TScert's proved verifier implements a monomorphic fragment $\Lang_0$
without effect variables, declared reach, $\kw{old}$, or capture-rooted
regions. Its regions are $\val$, $\st$, $\vw$, $\pi.\kw{spine}$,
$\pi.\kw{cells}$, $\pi.\kw{elemSpines}$, $\pi.f$, $\pi.\kw{fields}$, and
$\kw{all}$, rooted only at parameters and local variables. $\Lang_0$
splits what this paper unifies into two interpretations of a footprint:

\begin{center}\small
\begin{tabularx}{\textwidth}{@{}lX@{}}
\toprule
Verified fragment $\Lang_0$ (Lean) & Unified effect language\\
\midrule
Write footprint (\lean{Footprint}, \lean{Post}, \lean{Closed.Holds}) & Basic effect with $\kw{else}:\keep$\\
Fact support (\lean{Supported}) & Basic effect with $\kw{else}:\any$\\
Support soundness (\lean{Footprints.Supports}) & Fact stability under an effect (Def.~\ref{def:stable})\\
Frame respect (\lean{respects}, \lean{respect\_sound}, \lean{stable}) & Inclusion $\Phi\sqsubseteq_K\Psi$ with $\delta^\Psi_0=\any$\\
Closed callback effect & Variable-free effect\\
Regions (\lean{Footprints.Region}, \lean{Region.mem}) & Regions above\\
Change bounds (\lean{Change}, \lean{Change.Holds}, \lean{Change.le}) & Lattice $\Delta$\\
\bottomrule
\end{tabularx}
\end{center}

\noindent
In $\Lang_0$, Lean already proves (E1) by \lean{Region.mem}, (E3) by
\lean{Region.sub\_mem}, \lean{SepQuery.Sound}, and
\lean{Footprints.Region.apart\_of\_keys}, (E4) by
\lean{RigidityEvidence.sound}, and lattice monotonicity by
\lean{Change.Holds.mono}; (E2) is not stated as a separate Lean theorem.
Surface annotations in $\Lang_0$ cover \code{values}, \code{spines},
parameter spine regions (\code{"x[*]": "any"}), parameter field regions
(\code{"x.g"}, \code{"x.*"}, \code{"x.f.g"}, \code{"x.f.*"}),
\code{onThrow}, and \code{Fresh<T>}~\cite{ann}.

Because $\Lang_0$ lacks effect variables and capture roots, any
higher-order callback type must widen to a closed class effect (such as
\code{values: "any"; spines: "grow"}), which defeats W1 (cursor fields
in $\val$), W2 (handle fields in $\val$, asymmetric array growth in
$\st$), W3 (container writes colliding with the handle), and W4 (no way
to forward a client's effect).

\section{Prototype, Limits, and Related Work}\label{sec:status}

\subsection{Experimental Prototype}

To validate the rules on real Pretext code before refactoring TScert's
verified exit schemas, an experimental mode (\code{paper-rules-prototype})
layers these rules on top of TScert's ordinary pass:
\begin{itemize}[leftmargin=1.5em,itemsep=2pt]
\item \textbf{Architecture:} A source-level front-end pre-pass extracts
  effect bounds, instantiations, and \code{Effects<P, T>} signatures, and
  checks \ruleName{Inst}, module-local confinement, callback non-escape,
  and \code{Effects<P, T>} bodies coinductively (establishing cross-module
  interfaces via \code{--module-interface}). During the pass, an extension
  layer re-retains facts across polymorphic callback calls
  (\ruleName{Bound}, \ruleName{Anch}, \ruleName{Meet}), frames caller-side
  calls to quantified functions and monomorphic \code{Effects<P, T>}
  helpers (\ruleName{Confined}, \ruleName{Fields-Disj},
  \ruleName{Plain-Cells-Disj}, \ruleName{Elem-Len-Extend}), and injects
  scoped callback entry invariants (\ruleName{Callback-Inv}) for
  non-escaping inline closures. Outside non-pure variable calls, the layer
  agrees with the ordinary checker (proved, \lean{run\_eq\_base}). In
  addition, evaluating changes per location category
  (\lean{Change.onSpine}, \lean{Change.onValue}) provides a sound
  refinement of unprojected Definition~\ref{def:order}.
\item \textbf{Soundness boundary:} The prototype's verdict is distinct
  from TScert's verified verdict and rests on one named Lean axiom outside
  the \lean{Verification} library (\lean{paperKeptFactsSurvive}), together
  with the syntactic front-end checks (\ruleName{Inst}, confinement,
  \code{Effects<P, T>} bodies, non-escape) and scoped callback fact
  injection (\lean{closureFacts}). Neither \lean{mode8Accepted\_sound} nor
  any ordinary verdict depends on it.
\item \textbf{Tracked assumptions:} While verifying a library
  incrementally, missing source proofs or verifier gaps can be isolated
  with \code{paperAssume(fact, "\#N")}, checked as
  \code{if (!(fact)) throw "\#N"} in the target file only (never in
  imported modules, which are refused if they carry \code{paperAssume})
  and reported in the verdict. All 105/105 units of the Pretext evaluation
  (Section~\ref{sec:eval}) verify with zero \code{paperAssume} calls.
\end{itemize}

\subsection{Where the Prototype Conservatively Approximates the Rules}

The prototype never accepts more than the paper's rules, but decides
strictly less in the following ways (refusing where the rules would
accept):

\begin{center}\small
\begin{tabularx}{\textwidth}{@{}lX@{}}
\toprule
Rule / feature & Conservative approximation in \code{paper-rules-prototype}\\
\midrule
\ruleName{Anch} & Decided only for own fields and spines of locally allocated objects (from allocation separation facts), not nested elements or paths through fields.\\
Declared reach & Expands finite local and imported object, interface, array, typed-array, union, and intersection types (leaving generics, tuples, index signatures, and recursive types open); read at $\keep$ only, matched via pass read/alias links and reach inclusion ($\reach{\pi.f}{U}\subseteq\reach{\pi}{T}$).\\
Named effects & Named effect aliases (\code{type Geom<p> = ...}) are refused; bounds and \code{Effects<P, T>} require inline object type literals.\\
Region keys & Bounds and \code{Effects<P, T>} accept parameter-rooted paths \code{p.f$_1$.\dots.f$_k$[*]}, \code{p.f$_1$.\dots.f$_k$[\#]}, \code{p.f$_1$.\dots.f$_k$.f}, \code{p.f$_1$.\dots.f$_k$.*} ($k\ge 0$), \code{p.**}, and \code{"@state"}; \code{Effects<P, T>} additionally accepts class keys \code{fields} and \code{views} (no class keys in callback bounds; no structural region equality across distinct spellings). \ruleName{Plain-Cells-Disj} applies up to path depth~2, and \ruleName{Elem-Len-Extend} applies to linear lower-bound quantified clauses on element lengths.\\
\ruleName{Inst} \& anchoring & Disjointness and anchoring are decided syntactically before the pass over \code{const} locals, unassigned parameters, and \code{"@state"} (fresh literal locals are apart from parameter reaches when unescaped before the call; forwarded variables discharge callee bounds on own locations of caller-allocated fresh literals via \ruleName{Anch} and \ruleName{Meet}; paths, reaches, and cells are not decided anchored; \code{Join<E, P>} is checked componentwise).\\
\ruleName{Confined} & Checked locally per defining module for unexported module-level \code{let} variable cells (with interior writes through aliases charged to their target region, such as \code{fields}) and unexported \code{const} fresh containers (\code{Record}, \code{Map}, \code{WeakMap}, arrays, typed arrays) or lazy typed arrays that never escape; if loaded values are mutated, every stored value must be fresh and non-escaping.\\
\ruleName{Callback-Inv} & Non-escape is checked syntactically (each reference to the callback parameter must be a direct call, a nullish/truthiness test, or a direct forward to another non-escaping parameter), and entry facts are filtered using the base knowledge separation query alone (without type-disjoint plain-array cell separation or \ruleName{Elem-Len-Extend}).\\
Closures \& helpers & Callback types must carry a bare effect variable \code{Modifies<E, T>}. Closure effects are inferred from field/element writes, array \code{push}, fresh array/typed-array allocations, stated/forwarded calls, and pure/host natives; \code{Effects<P, T>} helper signatures are checked coinductively against their bodies outside Lean.\\
\ruleName{Bound} stability & Decided for single-location fields/spines at parameters, and for a reach only when the body's own effect (excluding callback calls) keeps the reach; calls through aliased callback locals or reassigned parameters are not read.\\
Call sites \& exits & Checks \ruleName{Inst} across function declarations in the project (refusing unanalysed value escapes of quantified functions); reads plain-variable, path, or pure-identity-wrapper arguments; refuses \code{onThrow: true} in effect-language literals (never erasing contracts); and over-approximates instantiated calls using their bounds.\\
\bottomrule
\end{tabularx}
\end{center}

\subsection{Limits of the Language}

Five open design limits remain beyond the prototype:
\begin{description}[leftmargin=1.5em,itemsep=2pt]
\item[Per-instance confined state] Module-private confined state
  ($\modstate{M}$) covers module-scoped caches such as W5; per-instance
  abstract regions ($x.@\kw{rep}$) for encapsulated classes and
  first-class maps remain open.
\item[Abstract exported effects] Exporting a named effect abstractly
  requires module-level obligations proving both its operation bounds and
  any advertised disjointness between handles.
\item[Explicit instantiations] Because TypeScript erases \code{Modifies}
  and \code{Effects}, call sites must write effect arguments explicitly.
\item[Exceptional exits] Bounding a callback on throws requires explicit
  \code{onThrow} annotations.
\item[Concurrency / \code{await}] Across an \code{await}, arbitrary
  interleaved code may run, so the effect widens to $\top$.
\end{description}

\subsection{Related Work}

\paragraph{Dynamic frames, region logic, higher-order framing, and monotonic state.}
Framing heap-dependent assertions without ownership types follows
\emph{dynamic frames}~\cite{kassios} and \emph{region logic}~\cite{bnr},
which track explicit footprints alongside state-dependent separation
judgments. Unlike standard dynamic frames and region logic---whose
per-location write sets are binary (unchanged vs.\ arbitrarily
modified)---our effects assign every region a relational bound from the
change lattice $\Delta$ (distinguishing $\keep$, $\init$, $\scalar$,
$\extend$, $\grow$, and $\any$). Monotonic-state
systems~\cite{pp,ahman} express similar growth properties via preorders
attached permanently to references at allocation time; here, bounds such
as $\extend$ and $\grow$ are per-step relational effect bounds, so the
same mutable array may grow across a callback in one phase and be shrunk
or frozen in another, while framing reduces to pointwise lattice
inclusion rather than arbitrary two-state theorem proving. Compared with
higher-order frame rules in separation logic~\cite{bty}, \ruleName{Anch}
encapsulates callee-private allocations purely relationally via the entry
heap mark $\kw{old}$ and anchoring, without ownership transfer or
separating conjunction.

\paragraph{Polymorphic effects and dependent types.}
Effect quantification and subeffecting follow Lucassen and
Gifford~\cite{lg} and Talpin and Jouvelot~\cite{tj}. In pure dependent
type theories, higher-order walkers avoid heap framing by threading
visitor state explicitly over immutable length-indexed vectors, while
imperative dependent systems such as Hoare Type Theory and F*'s
monotonic state~\cite{ahman} quantify over arbitrary two-state heap
predicates; our language trades arbitrary relational predicates for a
small finite change lattice whose inclusion is decidable on idiomatic
mutable code. Compared with inferring a callback's preservation demand
from a walker's body, writing a bounded effect variable
$\varepsilon\sqsubseteq\Phi$ (with $\delta^\Phi_0=\any$) makes the
contract explicit at the signature boundary, checks callers uniformly
via \ruleName{Inst}, preserves local cursors via \ruleName{Anch},
supports weaker bounds such as $\grow$, and propagates the callback's
effect into the walker's own signature.

\paragraph{TypeScript and JavaScript verification.}
Prior refinement and dependent type systems for this family---Dependent
JavaScript (DJS)~\cite{chugh} and Refined TypeScript (RSC)~\cite{vekris}
---verify array bounds and heap invariants by restricting mutable
aliasing (for example via immutability or uniqueness qualifiers), whereas
separation-logic toolchains such as JaVerT~\cite{javert} require explicit
separation assertions and footprints per routine. TScert's effect
language targets unannotated-ownership TypeScript directly:
type-driven declared reach ($\reach{\pi}{T}$) with key-based disjointness
and anchored effect variables handle mutable cursors, nested array
handles, and container-mutating callbacks without prohibiting aliasing.

\paragraph{Ownership and Rust verifiers.}
Ownership type systems such as Rust's enforce aliasing-XOR-mutability
at the language level, making W1, W2, and W4 largely free: a shared
borrow \code{\&PreparedLineData}, an unborrowed local cursor, and an
\code{FnMut} callback rule out handle and cursor mutation without
effect polymorphism, and disjoint closure captures cover W3 when
fields are accessed directly. Deductive and refinement verifiers for
Rust---including Prusti~\cite{prusti}, Creusot~\cite{creusot},
Flux~\cite{flux}, and Verus~\cite{verus}---leverage this ownership
discipline to eliminate most aliasing up front, focusing SMT reasoning
on arithmetic bounds, loop invariants, and \code{FnMut} contracts.
The trade-off is stylistic and structural: binary exclusivity
(\code{\&T} vs.\ \code{\&mut T}) forbids shared monotonic updates such
as $\grow$ across a live walker, breaks split borrows across
\code{\&mut self} helpers, and pushes transparent private caches (W5)
into interior mutability or explicit ghost permissions. TScert targets
the dual setting---idiomatic TypeScript with unrestricted aliasing---using
anchored effect polymorphism and the change lattice $\Delta$ to
recover framing without imposing an ownership discipline on source code.

\begin{thebibliography}{18}\small\itemsep1pt
\bibitem{tf} TScert Working Note. \emph{Tolerant frames: a modal account of fact preservation without ownership}. 27 November 2026.
\bibitem{patterns} TScert Working Note. \emph{Tolerant frames in practice: Pretext's patterns}. 27 November 2026.
\bibitem{fp} TScert Working Note. \emph{Footprints: stable operations for the facts constructs rely on}.
\bibitem{ann} TScert Documentation. \emph{Annotations}.
\bibitem{kassios} I.~Kassios. Dynamic frames: support for framing, dependencies and sharing without restrictions. FM 2006.
\bibitem{bnr} A.~Banerjee, D.~Naumann, S.~Rosenberg. Local reasoning for global invariants, part I: region logic. J.~ACM 60(3), 2013.
\bibitem{bty} L.~Birkedal, N.~Torp-Smith, H.~Yang. Semantics of separation-logic typing and higher-order frame rules. LICS 2005.
\bibitem{lg} J.~M.~Lucassen, D.~K.~Gifford. Polymorphic effect systems. POPL 1988.
\bibitem{tj} J.-P.~Talpin, P.~Jouvelot. Polymorphic type, region and effect inference. Journal of Functional Programming 2(3), 1992.
\bibitem{pp} A.~Pilkiewicz, F.~Pottier. The essence of monotonic state. TLDI 2011.
\bibitem{ahman} D.~Ahman, C.~Fournet, C.~Hri\c{t}cu, K.~Maillard, A.~Rastogi, N.~Swamy. Recalling a witness: foundations and applications of monotonic state. POPL 2018.
\bibitem{chugh} R.~Chugh, D.~Herman, R.~Jhala. Dependent types for JavaScript. OOPSLA 2012.
\bibitem{vekris} P.~Vekris, B.~Cosman, R.~Jhala. Refinement types for TypeScript. PLDI 2016.
\bibitem{javert} J.~Fragoso Santos, P.~Maksimovi\'c, D.~Naud\v{z}i\=unien\.e, T.~Wood, P.~Gardner. JaVerT: JavaScript verification toolchain. POPL 2018.
\bibitem{prusti} V.~Astrauskas, P.~M\"uller, F.~Poli, A.~J.~Summers. Leveraging Rust types for modular specification and verification. OOPSLA 2019.
\bibitem{creusot} X.~Denis, J.-H.~Jourdan, C.~March\'e. Creusot: a foundry for the deductive verification of Rust programs. ICFEM 2022.
\bibitem{flux} N.~Lehmann, A.~T.~Geller, N.~Vazou, R.~Jhala. Flux: liquid types for Rust. PLDI 2023.
\bibitem{verus} A.~Lattuada, T.~Hance, C.~Cho, M.~Brun, I.~Subasinghe, Y.~Zhou, J.~Howell, B.~Parno, C.~Hawblitzel. Verus: verifying Rust programs using linear ghost types. OOPSLA 2023.
\end{thebibliography}

\end{document}
